% ---------------------------------------------------------------------------
%  Section 2 -- Related Work   (~640 words)
%
%  Written independently of the MedSegLatDiff related-work section. Reusing
%  those sentences verbatim would be self-plagiarism even though the authors
%  overlap, so the two papers must not share prose.
%
%  The former 2.3 (Uncertainty Estimation and its Evaluation) was cut for
%  space: the MC dropout / ensemble contrast is already made in Sec. 3.4, and
%  the point about multi-annotator datasets is already made in Sec. 1. The
%  nguyen2025aleatoric citation was moved into 2.2 so it is not orphaned.
% ---------------------------------------------------------------------------

\section{Related Work}
\label{sec:related_work}

% ---------------------------------------------------------------------------
\subsection{Incomplete Multimodal Tumour Segmentation}
\label{subsec:rw-incomplete}

Brain tumour segmentation benchmarks assume four co-registered MRI sequences,
but clinical acquisitions often deviate. Examinations are shortened, series
are corrupted by motion, and contrast administration may be contraindicated. A
model that requires the complete protocol is therefore unusable on a large
share of real studies, and the literature on incomplete multimodal
segmentation addresses this in three main ways.

The first synthesises what is absent, generating the missing sequence and
passing the completed set to a conventional segmentation network. This leaves
the segmentation architecture unchanged, but accuracy then depends on the
quality of the synthesis, and errors introduced there carry through into the
mask. The second learns a representation that does not depend on which
sequences are present. HeMIS~\cite{hemis} pools per-modality features using
statistics that are defined for any subset, so a single model accepts
arbitrary inputs. RFNet~\cite{rfnet} refines this with region-aware fusion
that weights modalities differently for different tumour components, since
enhancement and oedema are visible in different sequences. The third distils a
full-modality teacher into a student that observes only a subset, transferring
knowledge the student cannot obtain from its own input. Transformer-based
approaches such as mmFormer~\cite{mmformer} combine modality-specific encoders
with inter-modality attention to model long-range dependencies across
sequences.

A recent refinement separates representation learning from fusion.
UniME~\cite{unime} first pretrains a single shared encoder with masked
self-supervision, so that a modality-agnostic representation is established
before any fusion mechanism is introduced. Parallel modality-specific encoders
are attached only afterwards, and their features are combined at multiple
scales. Decoupling the two stages reduces the burden on the fusion module,
which no longer has to learn what each sequence contributes and how to combine
them at the same time.

These methods differ in mechanism but agree in output: each returns a single
deterministic mask, whatever subset it was given. Accuracy is reported as a
function of the available sequences, but the confidence attached to the
prediction is not. A segmentation produced from one sequence is therefore
presented in the same form as one produced from four. How a model's expressed
confidence should respond when evidence is withdrawn is, to our knowledge, not
addressed in this literature.

% ---------------------------------------------------------------------------
\subsection{Generative Approaches to Segmentation}
\label{subsec:rw-generative}

Discriminative segmentation networks, from fully convolutional
architectures~\cite{fcn} through U-Net~\cite{unet} and its
successors~\cite{unet++,nnunet}, learn a mapping from image to mask. The
mapping is one-to-one by construction: for a given input the model has exactly
one answer, so the disagreement observed between expert annotators cannot
appear in the output.

Generative formulations relax this. Latent-variable models attach a stochastic
code to a segmentation backbone and decode different samples of that code into
different masks. The Probabilistic U-Net~\cite{probunet} draws a global latent
from a conditional prior, and PHiSeg~\cite{phiseg} extends the idea to a
hierarchy of latents at several resolutions. Their expressiveness is limited
by the size of that latent and by the tendency of the KL term to make the
decoder ignore it.

Diffusion models~\cite{ddpm} avoid this limit by building the distribution
through iterative refinement rather than a single stochastic bottleneck.
Conditioning a denoising network on the input image turns generation into
segmentation, and repeated sampling gives an ensemble from one set of
weights~\cite{ensemblediff,segdiff}. MedSegDiff~\cite{medsegdiff} strengthens
the conditioning with frequency-domain filtering of the guidance signal. The
cost is that every sample requires many network evaluations at full
resolution, which is acceptable in two dimensions and prohibitive in three.

Latent diffusion~\cite{ldm} removes that cost by moving the process into a
compressed space learned by an autoencoder. The approach has been used for
medical segmentation by LSegDiff~\cite{lsegdiff} and LDSeg~\cite{ldseg}.On the other hand, MedSegLatDiff~\cite{medseglatdiff} pairs two
autoencoders with a latent diffusion model and shows that a reconstruction
loss weighted towards the foreground is needed when targets are small, while
LatentFM~\cite{latentfm} replaces the diffusion trajectory with flow matching
in the same latent space to shorten sampling. Uncertainty arising from
incomplete inputs has also begun to receive attention, including flow-matching
formulations~\cite{nguyen2025aleatoric}. These methods operate on
two-dimensional images with all inputs present, and evaluate on datasets whose
ambiguity is fixed by the annotations. The present work extends this line to
volumetric data and to inputs that are deliberately incomplete.
